\documentclass[review,article]{planetary}
\addbibresource{references.bib}

\title{Mapping uncertainty in a small planetary-surface survey}
\shorttitle{Uncertainty in a planetary survey}
\author[1]{Maya Rivera}
\author[2]{Ifeoma Okafor}
\affil[1]{Example Planetary Science Centre, Example University}
\affil[2]{Remote Sensing Methods Lab, Example Institute}
\setcorrespauth{Maya Rivera}
\setcorrespmail{maya.rivera@example.edu}
\date{}

\begin{document}
\maketitle
\thispagestyle{firstpage}
\correspondingauthor
\submissiondate{29 August 2026}
\mstype{Research article}
\abstractwordcount{142}

\begin{abstract}
We demonstrate a transparent workflow for mapping uncertainty in a small
orbital-imaging survey. Three processing choices are varied independently, and
the resulting spread is reported beside the central estimate. The example is
intended as a practical starting point for a planetary-science manuscript.
\end{abstract}

\section{Introduction}
Planetary maps often combine measurements with processing choices that are
hard to compare after the fact. A useful manuscript therefore states which
inputs are fixed, which are varied, and how the final uncertainty is computed.
Here we use a synthetic crater-detection survey to make those choices explicit.

\section{Methods}
The survey contains 120 image tiles with a nominal ground sample distance of
12~m. A detection is retained when its confidence exceeds 0.80 and its fitted
diameter is between 0.5 and 4.0~km. We report the retained fraction $p$ and its
binomial standard error
\begin{equation}
  \sigma_p = \sqrt{\frac{p(1-p)}{n}},
  \label{eq:error}
\end{equation}
where $n$ is the number of inspected tiles. This makes the denominator visible
and avoids presenting a rounded percentage without its sample size.

\section{Results}
\begin{table}[ht]
\mytablestyle
\caption{Illustrative sensitivity analysis.}
\begin{tabular}{lrr}
\toprule
Processing choice & Retained tiles & Fraction \\
\midrule
Nominal threshold & 84 & 0.700 \\
Conservative threshold & 77 & 0.642 \\
Relaxed threshold & 91 & 0.758 \\
\bottomrule
\end{tabular}
\label{tab:sensitivity}
\end{table}

The nominal estimate is $p=0.700$, giving $\sigma_p\approx0.042$ from
Eq.~\eqref{eq:error}. The spread across processing choices is larger than the
sampling error, so it should be reported as a methodological sensitivity rather
than hidden inside a single confidence interval.

\section{Discussion}
The values in Table~\ref{tab:sensitivity} are a worked example, not a claim
about a real planetary body. Replace them with measurements, provide one file
per figure, and document the provenance of any external basemap or calibration
product. Persistent identifiers should be supplied for code and data whenever
available~\parencite{wilkinson2016}.

\section{Conclusion}
This starter covers the review layout, manuscript metadata, a methods equation,
a compact table, and an author-facing open-science reminder.

\printbibliography
\end{document}
